% Figure: partition of Z mod 4 --- four residue classes as bars.
% Source for Figure~\ref{fig:classes}. TikZ only; regenerable by compiling main.tex.
\begin{tikzpicture}[
  bar/.style={draw=rulegray,minimum height=7mm,font=\scriptsize,align=left},
  lab/.style={font=\footnotesize\sffamily,text=slate},
]
\node[bar,anchor=west,fill=accentsoft,text width=11.5cm] at (0,0)
  {$[0]_4=\{\dots,-8,-4,\mathbf{0},4,8,12,16,\dots\}$};
\node[bar,anchor=west,text width=11.5cm] at (0,-0.8)
  {$[1]_4=\{\dots,-7,-3,\mathbf{1},5,9,13,17,\dots\}$};
\node[bar,anchor=west,fill=accentsoft,text width=11.5cm] at (0,-1.6)
  {$[2]_4=\{\dots,-6,-2,\mathbf{2},6,10,14,18,\dots\}$};
\node[bar,anchor=west,text width=11.5cm] at (0,-2.4)
  {$[3]_4=\{\dots,-5,-1,\mathbf{3},7,11,15,19,\dots\}$};
\node[lab] at (6.2,-3.2) {four classes, zero overlap, covering all of $\Z$ ---
 $a\equiv b\pmod4$ iff they sit in the same bar};
\end{tikzpicture}
